\subsection{Complete local coefficient rings with finite residue field}
\label{C:complete-finite-residue-subsection}

We first supply an integral coefficient system. It will make
all finite levels below actual finite matrix problems, including
in mixed characteristic.

\begin{lemma}[Proved: explicit integral finite free coefficient coalgebras]
\label{C:integral-orthogonal-coalgebras}
Put $A=\mathbf Z[1/2]$ and let $O_m$ denote the split
orthogonal group. Its coordinate Hopf algebra $H_A=A[O_m]$
has an explicit increasing exhaustive filtration
$C_0\subset C_1\subset\cdots$ by finite free direct summands,
with explicitly computable $A$-bases and structural constants,
such that
\[
 \Delta(C_d)\subset C_d\otimes_A C_d,\qquad
 S(C_d)\subset C_d,\qquad
 C_d C_e\subset C_{d+e},\qquad 1\in C_0.              \tag{C.3.CR1}
\]
For $m\geq3$, $C_d$ is exactly the image of polynomials of
total degree at most $d$ in the matrix coordinates.
These statements, including the bases and inclusions,
remain true after arbitrary base change $A\to R$.
\end{lemma}
\begin{proof}
For $m\geq3$ use Cliff's orthogonal-standard bideterminants:
enumerate the pairs of tableaux of equal shape and size
at most $d$ satisfying the finite conditions of
\cite[Section 2, pp.\ 2--3]{C-CliffIntegralBasis2008}.
Their polynomials are products of the minors specified by
the two tableaux. Include the empty bideterminant $1$.
The straightening algorithm in
\cite[Theorem 5.1, p.\ 17; Section 7, proof of Theorem 7.1,
pp.\ 22--23]{C-CliffIntegralBasis2008}
uses coefficients in $A$ and never increases size.
The same proof establishes that all these standard
bideterminants form an $A$-basis of the integral coordinate
ring. We check precisely that this is the group-scheme
coordinate ring used here.

Let $J$ be the split symmetric Gram matrix in that source
and set $H=A[x_{ij}]/(X^tJX-J)$.
The relation implies $\det(X)^2=1$, so $X$ is invertible.
This algebra is smooth over $A$. Indeed, for a square-zero
ideal $I$ in an $A$-algebra, lift an orthogonal matrix
modulo $I$ to an invertible matrix $Y$. Write
$E=Y^tJY-J$, whose entries belong to $I$. Then
\[
                 Y'=Y(1-\tfrac12J^{-1}E)
\]
satisfies $Y'^tJY'=J$, because $E$ is symmetric and $I^2=0$.
Thus the algebra has the lifting property for square-zero
ideals. It is finitely presented, so is smooth and hence
flat over $A$
\cite[Proposition 10.138.13, Tag 00TN; Lemma 10.137.9,
Tag 00TA]{C-StacksCompleteCoefficients}.
In particular $H$ is torsion-free and embeds into $H\otimes_A
\mathbf Q$. The latter is the reduced coordinate ring
of the orthogonal group over $\mathbf Q$.
Its defining ideal is therefore the ideal of polynomials
vanishing on that group over an algebraic closure.
The defining ideal of $H$ is its contraction to $A[x_{ij}]$,
by torsion-freeness. This is precisely Cliff's integral
coordinate ring as defined in his Section 7.

It follows that the cited integral basis is a basis of
our quadratic quotient $H_A$ itself. Straightening and
the polynomial degrees prove the asserted description
of $C_d$. It is spanned by a subset of a basis, so both
$C_d$ and its complement are free. In particular all
inclusions remain injective after arbitrary base change.

Cauchy--Binet expands the coproduct of a minor as a sum
of tensor products of minors of the same size. For a
bideterminant of size $d$, it consequently gives
polynomials of degree at most $d$ in each tensor factor.
Straighten each factor to obtain the first containment
of (C.3.CR1) and its finite structural constants.
The antipode is induced by $X^{-1}=J^{-1}X^tJ$,
which is linear in the coordinates; it preserves degree.
Multiplication adds degrees. Straightening computes
the remaining constants as well. All these identities
are over $A$, so tensoring gives them over every $R$.

For completeness the small values of $m$ require no
tableau theorem. For $m=0$ take $C_d=A$ for every $d$.
For $m=1$ take
$C_d=A[z]/(z^2-1)$ with basis $1,z$ for every $d$,
and $\Delta(z)=z\otimes z$.
For $m=2$, the split group is
$\mathbf G_m\rtimes C_2$, where the nonidentity element
acts by inversion. This follows directly by writing
the diagonal matrices $\operatorname{diag}(u,u^{-1})$
and the corresponding antidiagonal matrices; the
quadratic equations give the two complementary
idempotents $ad$ and $bc$ in a matrix
$\left(\begin{smallmatrix}a&b\\c&d\end{smallmatrix}\right)$.
Write its coordinate ring as
$A[u,u^{-1}]e_0\oplus A[u,u^{-1}]e_1$.
Let $C_d$ have basis
\[
           u^j e_\epsilon,\qquad |j|\leq d,\quad
                            \epsilon\in\{0,1\}.
\]
Here $e_0+e_1=1$, and, with indices modulo two,
\[
 \Delta(u^j e_\eta)=
 \sum_{\epsilon+\delta=\eta}
 (u^j e_\epsilon)\otimes
 (u^{(-1)^\epsilon j}e_\delta).
\]
The antipode sends $u^j e_\epsilon$ to
$u^{-(-1)^\epsilon j}e_\epsilon$.
These formulas prove (C.3.CR1), exhaustiveness and
arbitrary base-change compatibility also in these cases.
\end{proof}

From now on $(R,\mathfrak m,\kappa)$ is a complete Noetherian
local ring with finite residue field of odd characteristic.
Put $R_N=R/\mathfrak m^N$, and use the fixed bases of
$C_d(R)=R\otimes_A C_d$ from
Lemma~\ref{C:integral-orthogonal-coalgebras}.
The rings $R_N$ are finite. Their multiplication tables
and transition maps will be called the \emph{finite quotient
access} to $R$; an individual element of $R$ is supplied
by its compatible residues. Statements about procedures
below always specify whether they need additional access.

For integers $n,d,N$, let $X_{n,d,N}$ be the finite set
of all $C_d(R_N)$-coactions on $R_N^n$.
Concretely these are the coefficient matrices satisfying
the counit equations and
$\Delta(\rho_{ij})=\sum_k\rho_{ik}\otimes\rho_{kj}$.
The fixed coalgebra basis turns these into finitely many
linear and quadratic equations in a finite ring.
The group $G_{n,N}=\operatorname{GL}_n(R_N)$ acts by
change of basis. Every set and action here is computed
by finite enumeration.

Fix total orders on the finite coefficient arrays modulo
$\mathfrak m^N$ so that each order refines the order of
its reduction at level $N-1$. Thus an entire fibre over
an earlier array precedes the entire fibre over a later
one. Order the elements within each fibre arbitrarily
but once and for all. Let $L_{n,d,N}$ be the finite list
of least elements of the $G_{n,N}$-orbits, found by
enumerating that finite group.

\begin{theorem}[Proved: explicit canonical inverse-limit module indices]
\label{C:complete-finite-canonical}
The sequences
\[
 (a_N)_{N\geq1},\qquad
 a_N\in L_{n,d,N},\quad a_{N+1}\bmod\mathfrak m^N=a_N,
                                                               \tag{C.3.CR2}
\]
are in bijection with isomorphism classes of rational
$O_{m,R}$-representations free of rank $n$ and with
coaction in $C_d(R)$.
Each sequence is an actual coaction on $R^n$.
No further isomorphism quotient is taken in (C.3.CR2).

An indexed module is indecomposable if and only if, at
some finite level $N$, its coefficient centralizer in
$M_n(R_N)$ contains no idempotent other than $0,I$.
Its least coefficient degree is $d$ if and only if,
when $d>0$, some finite level has a nonzero coefficient
outside $C_{d-1}(R_N)$.
Taking $n>0$ and imposing these two conditions gives
a repetition-free index set $\mathscr I_R$ for all
indecomposable finite projective rational representations.
Given an actual bounded coaction, every finite prefix
of its canonical sequence is computable from finite
quotient access.
\end{theorem}
\begin{proof}
Completeness of $R$ and freeness of the coefficient
coalgebra give
$X_{n,d}(R)=\varprojlim_N X_{n,d,N}$: coefficient
arrays lift uniquely, and each of the finitely many
polynomial identities holds exactly when it holds
modulo every power of $\mathfrak m$.
Likewise
$G_n(R)=\varprojlim_NG_{n,N}$.
The transition maps on these groups are surjective.
Indeed any lift of an invertible matrix is invertible
over a local ring because its determinant remains
nonzero in $\kappa$.

Consequently the reduction of the least element of
an orbit at level $N+1$ is the least element of its
orbit at level $N$: its orbit maps onto the entire
lower orbit, and the chosen orders refine reduction.
Thus reduction maps the lists $L_{n,d,N+1}$ to
$L_{n,d,N}$. A compatible sequence of these elements
is an actual coaction by the preceding paragraph.
Conversely, taking the least conjugate of an actual
coaction at every level gives such a sequence.

We check that the resulting coaction is isomorphic
to the original one and that distinct sequences are
not isomorphic. For two actual coactions $\rho,\rho'$,
let $Z_N$ be the set of $g\in G_n(R)$ intertwining
them modulo $\mathfrak m^N$. It is a closed subset
of the compact group $G_n(R)$, and the $Z_N$ are nested.
If the finite-level orbits agree, $Z_N$ is nonempty:
lift an element of $G_{n,N}$ witnessing this agreement.
Compactness gives an element in $\bigcap_N Z_N$.
The coefficient equations and separatedness of $R$
make it an actual intertwiner. The converse is immediate.
This proves the claimed bijection and uniqueness.

Let $\rho$ be an actual coaction.
Its decomposition as an $R$-projective representation
is equivalent to a commuting idempotent $e$ on $R^n$
other than $0,I$. Over a local ring a nonzero
idempotent has positive rank; hence a nontrivial one
has reduction modulo $\mathfrak m$ different from
both $0,I$. For each $N$, impose on $e\in M_n(R)$
the conditions $e^2=e$ and $e\rho=\rho e$ modulo
$\mathfrak m^N$, and the condition that its reduction
modulo $\mathfrak m$ is nontrivial.
These give nested compact sets. If a nontrivial
commuting idempotent exists at every finite level,
all these sets are nonempty and compactness gives
one over $R$. The converse follows by reduction.
This proves the finite-level indecomposability test.

The chosen coefficient filtration consists of basis
subsets. Membership of all coefficients in
$C_{d-1}(R)$ is therefore equivalent to the vanishing
of specified $R$-coordinates, which holds if and
only if it holds at every finite level. This proves
the assertion about least degree. Change of basis
preserves this degree in both directions, so retaining
its least value removes repetition between different
$d$. All finite projective modules over the local
ring $R$ are free, and every rational coaction uses
finitely many coordinate functions. Hence the
construction exhausts the category.
Finally each least finite-level conjugate is found
by an actual finite search, proving the assertion
about prefixes.
\end{proof}

For the next theorem, \emph{branch access} means the
following precise additional operation: given a node
in one of the finitely branching trees of finite
matrix solutions used below, decide whether it
extends to an infinite branch. A positive answer
also permits the branch to be generated by repeatedly
choosing the first child with a positive answer.
The trees will be given by explicit polynomial
matrix conditions over the finite rings $R_N$.
This is an access hypothesis, not a claim that
the operation is an ordinary computable one.

\begin{theorem}[Proved: explicit split form generators over complete finite-residue bases]
\label{C:complete-finite-WG}
The canonical index set $\mathscr I_R$ has a duality
and four explicitly specified subsets
$\mathscr I_{\rm s},\mathscr I_{\rm a},
\mathscr I_{\rm u},\mathscr I_{\rm h}$, defined in the
proof, for which
\[
 WG_{\rm sp}(BO_{m,R})\ \cong\
 \bigoplus_{i\in\mathscr I_{\rm s}}
       \bigl(\mathbf Z e_i\oplus(\mathbf Z/2)t_i\bigr)
 \ \oplus\
 \bigoplus_{i\in\mathscr I_{\rm a}\cup\mathscr I_{\rm u}
                         \cup\mathscr I_{\rm h}}
       \mathbf Z e_i.                                  \tag{C.3.CR3}
\]
The proof gives concrete representative forms for
every generator. With finite quotient access and
the specified branch access, it determines equality,
all product coefficients and all specified exterior
coefficients. Infinite coefficient arrays and indices
are output by their computable successive prefixes.
No uniform ordinary algorithm for equality of such
infinite data is asserted.
\end{theorem}
\begin{proof}
\emph{The complete endomorphism rings.}
For a finite free coaction $M$, the algebra
$E=\operatorname{End}_{O_{m,R}}(M)$ is the kernel of
a finite $R$-linear coefficient-commutator system
inside a finite matrix module. It is a finite
$R$-module because $R$ is Noetherian. It is complete
for the $\mathfrak m$-adic topology.
Moreover $\mathfrak mE\subset J=\operatorname{rad}E$:
for $x\in\mathfrak mE$, the geometric series for
$(1-x)^{-1}$ converges in $E$, and the same applies
after multiplication of $x$ by any element of $E$.
The finite-dimensional $\kappa$-algebra $E/\mathfrak mE$
has nilpotent radical. Therefore
$J^a\subset\mathfrak mE$ for some $a$, while
$\mathfrak mE\subset J$.
The $J$-adic and $\mathfrak m$-adic topologies agree,
so $E$ is $J$-adically complete.

Idempotents lift through $J$. One can successively
correct $e^2-e\in J$ by the convergent Newton formula
\[
          e\longmapsto e-(2e-1)^{-1}(e^2-e);
\]
the error is squared at each step.
Lifting primitive orthogonal idempotents successively
in complementary corners gives a complete set with
local corner rings. Their images on $M$ are finite
projective $R$-modules and hence free.
Thus the category has Krull--Schmidt decompositions,
split idempotents, and endomorphism rings complete
in their radicals. These are precisely the hypotheses
of the complete-radical form theorem of
\cite[Theorem 2.2, Corollary 2.3, Propositions 2.4--2.5,
Theorems 3.2--3.3, pp.\ 269--273]{C-QSS1979}.
In particular hermitian cancellation holds.

\emph{The residue data without an unknown module index.}
Let $M_i$ be the actual canonical coaction given by
an index in $\mathscr I_R$. Its dual is obtained
by the antipode and transpose, and then by the
canonical construction of
Theorem~\ref{C:complete-finite-canonical}.
Its endomorphism ring $E_i$ is local. Its residue
division ring $D_i=E_i/\operatorname{rad}E_i$ is finite:
it is a quotient of the finite algebra
$E_i/\mathfrak mE_i$.
It is consequently a finite field, by
\cite[Theorem III.4.4, p.\ 29]{C-ArtinFiniteDivision}.

These residue fields can be computed using exactly
the stated access. Let $B_i$ be the image of $E_i$
in $\operatorname{End}_\kappa(M_i/\mathfrak mM_i)$.
For every matrix in this finite ambient set, the tree
of lifts satisfying the coefficient-commutator
equations decides membership in $B_i$ by branch access.
The kernel of $E_i\to B_i$ lies in its radical:
an endomorphism with matrix entries in $\mathfrak m$
has $1-uh$ invertible for every $u\in E_i$, since
its determinant is $1$ modulo $\mathfrak m$.
Its matrix inverse is again equivariant.
Thus
$D_i=B_i/\operatorname{rad}B_i$.
This quotient and its finite field operations are
found by finite enumeration in $B_i$.

If $M_i$ is self-dual, choose a nonsingular invariant
symmetric form if one exists, and otherwise an
alternating form. One exists by the local-ring
argument $2g=(g+g^*)+(g-g^*)$ and the cited
Proposition 2.5. Choose the least branch of the
corresponding finite matrix solution tree, with
determinant nonzero modulo $\mathfrak m$.
Call the result $q_i$, with sign $\varepsilon_i$.
Its adjoint involution on $E_i$ induces an explicitly
computed involution $\sigma_i$ on $D_i$.
Define $\mathscr I_{\rm s}$ by
$\varepsilon_i=1,\sigma_i=1$,
$\mathscr I_{\rm a}$ by
$\varepsilon_i=-1,\sigma_i=1$, and
$\mathscr I_{\rm u}$ by $\sigma_i\ne1$.
Let $\mathscr I_{\rm h}$ contain the smaller canonical
index in each non-self-dual pair. Distinct infinite
canonical arrays are ordered by their first different
entry, after ordering their finite rank and degree.
Equality and the requisite choices use branch access
to the explicit intertwiner trees.

\emph{Forms and generators.}
Complete-radical reduction and transfer identify the
self-dual block with $GW^{\varepsilon_i}(D_i,\sigma_i)$.
The finite-field classification
\cite[Lemma 3.1(a),(b)]{C-SergeichukFiniteForms}
gives $\mathbf Z\oplus\mathbf Z/2$ in symmetric
trivial-involution type, and $\mathbf Z$ in alternating
or nontrivial-involution type. Non-self-dual pairs
give the hyperbolic $\mathbf Z$ summands.
This proves (C.3.CR3), including its split rather
than Witt convention.

Here are actual representatives. In symmetric type
put $e_i=[q_i]$. Choose the first nonsquare
$c\in D_i$, lift it to $b\in E_i$ by the least
endomorphism branch with its prescribed residue,
and set $a=(b+b^*)/2$.
Then $a$ is an invertible self-adjoint endomorphism.
Put $t_i=[q_i a]-[q_i]$; its exact order is two
because its residue determinant is a nonsquare.
For alternating type and for a non-self-dual pair
put $e_i=H(M_i)$.
For unitary type, choose a nonzero $c\in D_i$ with
$\sigma_i(c)=\varepsilon_i c$, lift it to $b$, and
put $a=(b+\varepsilon_i b^*)/2$.
Then $q_i a$ is symmetric and nonsingular; its
class is $e_i$. Such $c$ exists: use $1$ when
$\varepsilon_i=1$, and $u-\sigma_i(u)$ otherwise.
The complete-radical theorem proves that the displayed
representatives have exactly the stated generators
and relations.

\emph{Operations and their exact access.}
To split a given module, use the finite trees of
nontrivial commuting idempotents. Branch access
decides whether a split exists and supplies it.
Choose a basis for its image by lifting a nonzero
minor of its reduction modulo $\mathfrak m$.
The idempotent makes the image projective; Nakayama
and freeness give the resulting basis over $R$.
Repeat on both summands. Their positive ranks
decrease, so at most $n-1$ splits occur.
For each summand, branch access decides which
of its finitely many coefficient coordinates
are zero, and hence its least degree; its canonical
prefixes are then given by finite orbit enumeration.

For a resulting actual symmetric form, the underlying
Krull--Schmidt multiplicities leave finitely many
possible form coordinates: one determinant bit
for each occurring symmetric type, and only
the multiplicity in each other type. More explicitly,
a symmetric multiplicity $r>0$ and bit $b$ is represented
by $r-b$ copies of $q_i$ and $b$ copies of $q_i a$;
its coordinate is $r e_i+b t_i$.
An alternating multiplicity is even and contributes
half its multiplicity; a unitary multiplicity
contributes itself; a paired block contributes
the common multiplicity of its two dual constituents.
Construct the finitely many candidate forms and
test their equivariant isometry by branch access
to the corresponding finite matrix trees.
The complete-radical classification shows that
exactly one coordinate choice succeeds.

Tensor coactions are products of coefficients and
exterior coactions are their minors. Their bounds
are $d+e$ and $rd$, respectively, by (C.3.CR1).
Tensor Gram matrices and exterior determinant forms
are calculated over $R$; no factorial is inverted.
Applying the preceding decomposition gives every
structure coefficient in (C.3.CR3).
Bilinearity and
$\lambda_t(-x)=\lambda_t(x)^{-1}$ give arbitrary
virtual inputs, to every specified finite order.
All searches except the stated branch decisions
are finite; matrices or indices with infinitely
many digits are generated one prefix at a time.
This proves the precise access assertion without
claiming a finite ordinary equality test.
\end{proof}

\begin{theorem}[Proved: an exact bounded-image description using solved finite rings]
\label{C:complete-finite-bounded-image}
Put
$G_N=WG_{\rm sp}(BO_{m,R_N})$, with the explicit
canonical coordinates and transition algorithms of
Theorem~\ref{C:Om-finite-base-algorithm}.
Reduction defines an injective homomorphism, compatible
with every exterior operation,
\[
 WG_{\rm sp}(BO_{m,R})\ \longrightarrow\
                       \varprojlim_N G_N.                \tag{C.3.CR4}
\]
Its image consists exactly of the compatible sequences
$(\xi_N)$ with the following explicit property:
there are integers $B,D\geq0$ such that, for every $N$,
$\xi_N$ is a difference of two actual nonsingular
symmetric rational representations, each of rank
at most $B$ and with coefficients in $C_D(R_N)$.

For fixed $B,D,N$ this is a finite decidable test,
using only the finite ring $R_N$ and its solved
canonical form coordinates. Consequently this gives
another description of the entire split ring as
a literal set of compatible sequences subject to
an explicit $\exists B,D\,\forall N$ finite predicate,
with no remaining isomorphism quotient.
Products and specified exterior coefficients are
computed coordinatewise by finite procedures;
their successive prefixes require only finite
quotient access. Equality of infinite sequences
is not asserted to be decided by finitely many
of those computations.
\end{theorem}
\begin{proof}
Every element of the split group is $[P,q]-[Q,h]$
for two actual symmetric representations. They
have finite ranks and finite coefficient bounds,
so their reductions satisfy the stated uniform
boundedness condition.
If this difference is zero at every level,
hermitian cancellation over the finite local ring
$R_N$ makes the two reduced forms actually
equivariantly isometric at every level.
Cancellation follows either from the finite-base
form classification or from
\cite[Section 3.4(1), p.\ 273]{C-QSS1979}.
The two ranks are equal, already at the first level.
For each $N$ consider all matrices in
$\operatorname{GL}_n(R)$ which are equivariant
isometries modulo $\mathfrak m^N$.
They form nonempty nested closed subsets of a
compact group. Their intersection is nonempty
and, by the finite coefficient equations and
separatedness, consists of actual equivariant
isometries. Thus the original difference is zero.
This proves injectivity.

Conversely suppose a compatible sequence $(\xi_N)$
has bounds $B,D$. At level $N$ enumerate all pairs
of coactions in $C_D(R_N)$ of ranks at most $B$,
and all symmetric Gram matrices with determinant
a unit. Keep precisely the pairs for which both
forms are invariant and their difference has
canonical finite-base coordinate $\xi_N$.
This is a finite set, and the boundedness assumption
says it is nonempty. Reduction sends this set
at level $N+1$ into the corresponding set at
level $N$, because the $\xi_N$ are compatible.
These finite sets therefore form a finitely
branching tree with nonempty levels of arbitrary
depth. Such a tree has an infinite branch:
choose a first-level vertex with descendants
of arbitrarily large depth, then a child with
the same property, and continue.

Along this branch the two ranks are constant.
The compatible coefficient arrays lift to coactions
on free $R$-modules by completeness and the finite
free coalgebra equations. The compatible Gram
matrices lift as well. They are symmetric and
invariant, and their determinants are units
because they are nonzero modulo $\mathfrak m$.
Thus the branch gives actual forms $(P,q),(Q,h)$
over $R$, whose difference maps to $(\xi_N)$.
This proves the exact image assertion.

The finite test in the theorem is precisely the
enumeration just described; equality of its
candidate coordinates with $\xi_N$ is the already
proved finite-base equality algorithm.
All transition maps are likewise computed by
reducing actual generator coactions and Gram
matrices and using that algorithm.
There is therefore no unclassified finite-level
quotient hidden in (C.3.CR4).

Reduction preserves tensor products and exterior
forms, so all operations in (C.3.CR4) are coordinatewise.
The image condition is preserved by them.
For example, if $x=P-Q$ and $y=P'-Q'$ have rank
bounds $B,B'$ and coefficient bounds $D,D'$,
then
\[
 xy=(P\otimes P'\oplus Q\otimes Q')
       -(P\otimes Q'\oplus Q\otimes P')
\]
has rank bound $2BB'$ and coefficient bound $D+D'$.
For a fixed exterior coefficient $r$,
expand $\lambda_t(P)\lambda_t(Q)^{-1}$ to order $r$.
This is a finite integral combination of tensor
products of exterior powers, with a rank bound
depending only on $B,r$ and coefficient degree
at most $rD$. Group its positive and negative
terms to obtain the required two actual forms.
The finite-base algorithms compute each resulting
coordinate. This proves all remaining assertions.
\end{proof}

\begin{theorem}[Proved: arbitrary products of the complete finite-residue bases]
\label{C:complete-finite-products}
Let $(R_i,\mathfrak m_i)_{i\in I}$ be any family of complete
Noetherian local rings with finite residue fields of odd
characteristic, and put $R=\prod_{i\in I}R_i$.
No Noetherian hypothesis is imposed on this product.
There is an injective map, preserving products and all
exterior operations,
\[
 WG_{\rm sp}(BO_{m,R})\longrightarrow
                \prod_{i\in I}WG_{\rm sp}(BO_{m,R_i}).    \tag{C.3.CR5}
\]
Its image consists exactly of the families admitting
one uniform rank bound $B$ and one uniform coefficient
bound $D$: each coordinate is a difference of two
actual symmetric forms of ranks at most $B$ and
coaction degree at most $D$.

Equivalently, the entire group and its operations have
the following explicit description using only solved
finite rings. Take the arrays
\[
  (\xi_{i,N})_{i\in I,\ N\geq1},\qquad
  \xi_{i,N}\in WG_{\rm sp}
        (BO_{m,R_i/\mathfrak m_i^N}),
\]
compatible in $N$ for each $i$, and impose the single
uniform condition
\[
 \begin{gathered}
 \exists B,D\quad\forall i\in I\quad\forall N\geq1:\\
 \xi_{i,N}\text{ has a pair of representing forms}\\
 \text{of ranks }\leq B\text{ and degrees }\leq D.
 \end{gathered}                                       \tag{C.3.CR6}
\]
Each fixed-coordinate test is the finite decidable
test of Theorem~\ref{C:complete-finite-bounded-image}.
The operations are coordinatewise, and each specified
coordinate of each operation is computed by the
finite-base algorithms. The assertion does not
claim a finite procedure for the universal
quantifiers over $I$ or all $N$.
\end{theorem}
\begin{proof}
A finite projective $R$-module is the image of an
idempotent in $M_b(R)$ for some finite $b$.
Its specialization to each local factor is free
of rank at most $b$.
A rational coaction on it is represented, using
this idempotent, by finitely many elements of
$R[O_m]$. They belong to $C_D(R)$ for some $D$.
Changing the bases over the local factors uses
constant matrices and cannot increase this bound.
Thus a difference of two global actual forms
has the asserted uniform bounds.

Conversely suppose the local coordinates are
represented by forms $(P_i,q_i)-(Q_i,h_i)$ with
the stated uniform bounds. Choose bases
$P_i\simeq R_i^{r_i}$, where $r_i\leq B$.
Let $e_i\in M_B(R_i)$ be the diagonal idempotent
whose first $r_i$ entries are $1$ and the rest $0$.
The family gives an idempotent
$e=(e_i)\in M_B(R)$ and a finite projective module
$P=eR^B$.
Extend the coaction on $P_i$ to $R_i^B$ by giving
the complementary coordinate summand the trivial
coaction. All its coefficients lie in $C_D(R_i)$,
since $1\in C_0$.
The fixed finite free basis of $C_D$ allows these
coefficient matrices to be assembled entry by
entry into matrices over $R$.
The counit and coassociativity identities hold
coordinatewise, hence over $R$; the same is true
of equivariance of $e$.
Restricting gives the required coaction on $P$.

Extend the Gram matrix $q_i$ and its inverse
by zero on the complement. Their entries likewise
assemble over $R$, and their inverse identities
on the images of the $e_i$ assemble to the inverse
identities on $P$. Symmetry and equivariance are
coordinatewise identities, so they give a
nonsingular invariant symmetric form $q$ on $P$.
Perform the same construction for the $Q_i$.
The resulting difference of global forms has
all the prescribed coordinates. This proves
the stated image.

For injectivity write a global element as
$[P,q]-[Q,h]$ and suppose it becomes zero in
every local factor. Complete-radical hermitian
cancellation, verified in
Theorem~\ref{C:complete-finite-WG}, gives an actual
equivariant isometry $P_i\simeq Q_i$ for every $i$.
Choose finite idempotent presentations of $P,Q$.
Composing the local isometries and their inverses
with the inclusions and projections of these
presentations gives matrices of two fixed finite
sizes over each $R_i$. Assemble their entries
over $R$. Their inverse identities and Gram
identities hold over $R$.
Their equivariance equations lie in one finite
free coefficient coalgebra and hold in every
coordinate; they therefore hold over $R$.
Thus the assembled maps are an actual global
equivariant isometry, proving injectivity.

Finally apply
Theorem~\ref{C:complete-finite-bounded-image}
to each local factor. If (C.3.CR6) holds, that
theorem supplies the two local actual forms
with the same fixed bounds $B,D$, so the assembly
just proved applies. The converse was already
shown. This gives the exact array description,
without an unresolved representation-isomorphism
quotient.
Tensor products and exterior forms commute with
base change, and the explicit uniform bounds
in the proof of
Theorem~\ref{C:complete-finite-bounded-image}
remain uniform over all $i$.
Hence every operation is the stated coordinatewise
operation on this exact image.
\end{proof}
